% Additional lessons: arrays, bits, graphs, lists, and strings.

\patternintoc{Two Pointers}
\problemstart{Container With Most Water}{11}{Medium}{Two pointers}{Array}

\subsection{The Problem}
Each array value is the height of a vertical line. Choose two lines that form
the container with the greatest area. For indexes \(l<r\): \(\text{area}=\min(h_l,h_r)(r-l).\)

\begin{conceptbox}{Example}
For \texttt{[1,8,6,2,5,4,8,3,7]}, lines at indexes \(1\) and \(8\) create
area \(\min(8,7)\cdot7=49\).
\end{conceptbox}

\subsection{From Brute Force to Two Pointers}
Checking every pair takes \(O(n^2)\). Start with the widest possible pair.
The shorter line limits the area. Moving the taller line inward only reduces
width while keeping the same limiting height or worse, so it cannot improve
that pair. Move the shorter line and search for a taller boundary.

\begin{invariantbox}{every removed wall is safely dominated}
Before each iteration, every pair outside the current interval has already
been measured or proved unable to beat a measured pair. Once the current area
is recorded, any narrower pair that keeps the shorter wall has no greater
limiting height, so that wall may be discarded.
\end{invariantbox}

\begin{visualbox}{The shorter wall limits how much water fits}
\centering
\begin{tikzpicture}[>={Stealth[length=2.2mm,width=1.55mm]},font=\small,x=8mm,y=4mm]
  \foreach \x/\h in {0/1,1/8,2/6,3/2,4/5,5/4,6/8,7/3,8/7} {
    \draw[navy,very thick] (\x,0)--(\x,\h);
    \node[below] at (\x,0) {\scriptsize \x};
  }
  \begin{scope}[on background layer]
    \fill[softteal,opacity=.65] (1,0) rectangle (8,7);
  \end{scope}
  \draw[teal,very thick] (1,0)--(1,8); \draw[orange,very thick] (8,0)--(8,7);
  \node[above=2mm,text=teal,font=\bfseries] at (1,8) {left: 8};
  \node[above=2mm,text=orange,font=\bfseries] at (8,7) {right: 7};
  \draw[decorate,decoration={brace,amplitude=4pt},navy]
    (1,-1.1)--(8,-1.1) node[midway,below=4mm]{width \(=7\)};
  \node[fill=white,rounded corners,inner sep=2mm] at (4.5,3.5)
    {area \(=7\times7=49\)};
\end{tikzpicture}
\end{visualbox}

\subsection{Pseudocode and C++ Implementation}
\begin{lstlisting}[style=pseudocode]
left = 0, right = n - 1, best = 0
while left < right:
    best = max(best, min(height[left], height[right]) * (right-left))
    if height[left] <= height[right]: left++
    else: right--
return best
\end{lstlisting}

\begin{lstlisting}[style=compactcode,caption={Complete C++ solution: Container With Most Water}]
#include <algorithm>
#include <vector>

using namespace std;

class Solution {
public:
    int maxArea(vector<int>& height) {
        int left = 0;
        int right = static_cast<int>(height.size()) - 1;
        long long best = 0;

        while (left < right) {
            const long long width = right - left;
            const long long wall = min(height[left], height[right]);
            best = max(best, width * wall);
            if (height[left] <= height[right]) ++left;
            else --right;
        }
        return static_cast<int>(best);
    }
};
\end{lstlisting}

\subsection{Why It Works, Time, and Space}
When a pair is processed, every pair that keeps its shorter wall and uses a
smaller width is no better. Discarding that wall cannot remove the best
unseen pair. Each pointer moves at most \(n-1\) times, so time is \(O(n)\) and
space is \(O(1)\).

\begin{revisionbox}{Common mistake}
Area uses the smaller height, not the larger one. This is not binary search:
the pointers move according to which boundary limits the current container.
\end{revisionbox}

\subsection{Detailed Pointer Dry Run}
\begin{dryrunbox}{Worked example \texttt{[1,8,6,2,5,4,8,3,7]}}
\begin{center}
\begin{tabularx}{\linewidth}{@{}c c c c X@{}}
\toprule
Left, right & Heights & Width & Area & Pointer removed \\
\midrule
0, 8 & 1, 7 & 8 & 8 & left: height 1 limits the pair \\
1, 8 & 8, 7 & 7 & 49 & right: height 7 is shorter \\
1, 7 & 8, 3 & 6 & 18 & right \\
1, 6 & 8, 8 & 5 & 40 & either equal wall may move \\
2, 6 & 6, 8 & 4 & 24 & left \\
\bottomrule
\end{tabularx}
\end{center}
No later pair exceeds 49, so the answer is the container between indexes 1
and 8.
\end{dryrunbox}

\subsection{Why Moving the Shorter Wall Is Safe}
Suppose the left wall is shorter. Keeping it while moving the right pointer
inward reduces the width and cannot increase the limiting height above that
same left wall. Every such pair is no better than the current pair. The only
chance to improve is to discard the shorter wall and search for a taller one.
This reason is the key rule behind the two-pointer method.

\subsection{Code Walkthrough}
The code begins with the widest possible pair. It calculates width and the
smaller boundary with \texttt{long long} intermediates, updates the best area,
and advances exactly one pointer. Since every iteration shrinks the interval,
the loop terminates after at most \(n-1\) moves.

\subsection{Edge Cases and Follow-ups}
Two lines form the smallest valid input. Zero-height lines are allowed and
simply create zero area. Equal boundary heights may move either side because
the current limiting height cannot improve while the width shrinks. A common
follow-up asks why sorting is impossible: the original index distance is part
of the area and must be keepd.

\begin{revisionbox}{How to explain it simply}
I start with maximum width. The shorter wall limits the current area, and any
narrower pair that keeps that wall cannot do better. I move the
shorter side, test each remaining choice once, and obtain \(O(n)\) time and
\(O(1)\) space.
\end{revisionbox}
\chaptercheckpoint
  {Two boundaries determine an area, and moving inward always reduces width.}
  {Left and right pointers plus the best area seen so far.}
  {Moving the taller wall without proving which boundary is safely discardable.}
  {$O(n)$ time and $O(1)$ extra space.}
\judgeclassnote


\patternintoc{Arrays, Hashing, and Bits}
\problemstart{Contains Duplicate}{217}{Easy}{Membership hashing}{Hash set}

\subsection{The Problem}
Return true when any integer appears at least twice. Return false when every
value is distinct.

\subsection{How We Get to the Solution}
Comparing every pair costs \(O(n^2)\). Sorting costs \(O(n\log n)\) and changes
the input order. A hash set answers the exact question we need: have we seen
this value before?

\begin{invariantbox}{the set equals the distinct processed prefix}
Immediately before reading \texttt{nums[i]}, the set contains every distinct
value at an earlier index and no future value. A membership hit is a
real duplicate; otherwise insertion keeps the statement for the next
iteration.
\end{invariantbox}

\begin{visualbox}{Stop when a number appears again}
\centering
\begin{tikzpicture}[>={Stealth[length=2.2mm,width=1.55mm]},font=\small]
  \foreach \x/\v in {0/1,1/2,2/3,3/1} {
    \node[arraycell] (a\x) at (1.2*\x,0) {\v};
  }
  \node[dsnode,right=22mm of a3] (s) {numbers already seen\\\{1,2,3\}};
  \draw[flowarrow,orange] (a3)--node[edgelabel,above,font=\scriptsize]{seen 1 before} (s);
  \node[below=7mm of s,draw=orange,fill=softorange,rounded corners] {duplicate found};
\end{tikzpicture}
\end{visualbox}

\subsection{Pseudocode}
\begin{lstlisting}[style=pseudocode]
seen = empty set
for value in nums:
    if value is already in seen: return true
    insert value into seen
return false
\end{lstlisting}

\subsection{C++ Implementation}
\begin{lstlisting}[style=compactcode,caption={Complete C++ solution: Contains Duplicate}]
#include <unordered_set>
#include <vector>

using namespace std;

class Solution {
public:
    bool containsDuplicate(vector<int>& nums) {
        unordered_set<int> seen;
        seen.reserve(nums.size());
        for (const int value : nums) {
            if (!seen.insert(value).second) return true;
        }
        return false;
    }
};
\end{lstlisting}

\subsection{Complexity and Revision}
Hash insertion is average-case \(O(1)\), giving average \(O(n)\) time and
\(O(n)\) space. The worst-case hash bound is \(O(n^2)\), although interview
analysis normally states the average case. Empty and one-element arrays return
false.

\begin{revisionbox}{Pattern clue}
Questions asking whether an item has appeared before usually suggest a set.
Questions needing information attached to that item suggest a map.
\end{revisionbox}

\subsection{Example Walkthrough}
For \texttt{[1,2,3,1]}, the set evolves as follows:
\begin{dryrunbox}{Insertion result reveals the duplicate}
\begin{center}
\begin{tabular}{@{}c c c@{}}
\toprule
Current value & Set before insertion & Result \\
\midrule
1 & \(\{\}\) & insert succeeds \\
2 & \(\{1\}\) & insert succeeds \\
3 & \(\{1,2\}\) & insert succeeds \\
1 & \(\{1,2,3\}\) & insert fails; return true \\
\bottomrule
\end{tabular}
\end{center}
\end{dryrunbox}

The set is enough because only membership matters. A map would store an
unneeded second value.

\subsection{Code Walkthrough and Alternatives}
\texttt{seen.reserve(nums.size())} reduces avoidable rehashing but is not
required for why the solution works. \texttt{insert} returns a pair whose Boolean field
is false when the key already exists, allowing lookup and insertion in one
operation. Returning immediately avoids scanning the rest of the array.

Sorting is an alternative: sort, then compare adjacent values in
\(O(n\log n)\) time and constant or implementation-dependent extra space.
The nested-loop baseline uses \(O(1)\) space but \(O(n^2)\) time.

\subsection{Edge Cases and Interview Notes}
Empty and one-element inputs are distinct by definition. Negative values and
zero require no special handling. Repeated values may appear more than twice;
the second occurrence is enough. Hash operations are average-case constant
time, not a deterministic worst-case guarantee.

\begin{revisionbox}{Practice checkpoint}
Change the question to ``return the first duplicated value'' and explain why
the same scan works. Then change it to ``return each value's frequency'' and
identify why a map is now required.
\end{revisionbox}
\chaptercheckpoint
  {The question asks whether a value has appeared earlier.}
  {A hash set containing the distinct processed values.}
  {Forgetting that hash bounds are average case or using a map when only membership matters.}
  {$O(n)$ average time and $O(n)$ space.}
\judgeclassnote


\problemstart{Counting Bits}{338}{Easy}{Bit dynamic programming}{Array}

\subsection{The Problem}
For every integer from \(0\) through \(n\), return the number of set bits in
its binary representation. For \(n=5\), return \texttt{[0,1,1,2,1,2]}.

\subsection{Update Rule}
Counting every number bit by bit takes \(O(n\log n)\). Dividing \(i\) by two
removes its last binary digit, and \(i\&1\) tells us whether that removed digit
was one: \(bits[i]=bits[i\!>>\!1]+(i\&1).\)

\begin{invariantbox}{the shifted dependency is already solved}
When the loop reaches \(i\), the index \(i>>1\) is smaller than \(i\), so its
bit count is already correct. Adding the removed least-significant bit gives
the exact count for \(i\) and extends the solved prefix by one entry.
\end{invariantbox}

\begin{visualbox}{Use the answer for 2 to find the answer for 5}
\centering
\begin{tikzpicture}[font=\small,>={Stealth[length=2.2mm,width=1.55mm]},node distance=10mm]
  \node[dsnode,active] (five) {\(5=101_2\)};
  \node[dsnode,right=of five] (shift) {remove last digit\\\(101\to10\)};
  \node[dsnode,right=of shift,visited] (known) {ones in 2\\\(=1\)};
  \node[dsnode,below=12mm of shift,fill=softorange] (last) {removed digit\\is 1};
  \node[dsnode,right=14mm of last,active] (answer) {ones in 5\\\(=2\)};
  \draw[flowarrow] (five)--(shift);
  \draw[flowarrow] (shift)--(known);
  \draw[flowarrow,teal] (known.south) to[bend right=18] (answer.north);
  \draw[flowarrow,orange] (last)--(answer);
  \foreach \x/\v in {0/0,1/1,2/1,3/2,4/1,5/2} {
    \node[arraycell,minimum width=10mm] (bit\x) at ({-0.2+1.12*\x},-3.35) {\v};
    \node[below=1mm of bit\x,font=\scriptsize] {\x};
  }
  \node[left=8mm of bit0,text=navy,font=\bfseries,align=right] {answers for\\0 through 5};
  \node[draw=orange,very thick,fit=(bit5),inner sep=1mm] {};
\end{tikzpicture}
\end{visualbox}

\begin{dryrunbox}{Binary update rule}
\begin{center}
\begin{tabular}{c|cccccc}
\(i\) & 0 & 1 & 2 & 3 & 4 & 5 \\ \hline
binary & 0 & 1 & 10 & 11 & 100 & 101 \\
set bits & 0 & 1 & 1 & 2 & 1 & 2
\end{tabular}
\end{center}
\end{dryrunbox}

\subsection{Pseudocode}
\begin{lstlisting}[style=pseudocode]
bits[0] = 0
for value from 1 through n:
    bits[value] = bits[value / 2] + (value mod 2)
return bits
\end{lstlisting}

\subsection{C++ Implementation}
\begin{lstlisting}[style=compactcode,caption={Complete C++ solution: Counting Bits}]
#include <vector>

using namespace std;

class Solution {
public:
    vector<int> countBits(int n) {
        vector<int> bits(n + 1, 0);
        for (int value = 1; value <= n; ++value) {
            bits[value] = bits[value >> 1] + (value & 1);
        }
        return bits;
    }
};
\end{lstlisting}

\subsection{Why It Works and Complexity}
The update rule reuses the already-computed answer for \(\lfloor i/2\rfloor\)
and adds the final bit. Time is \(O(n)\); extra working space is \(O(1)\)
excluding the required \(O(n)\) output. The base \(bits[0]=0\) is essential.

\subsection{Building the Bit Update Rule}
Every non-negative integer \(i\) can be written as \(i=2\left\lfloor i/2\right\rfloor +(i\bmod 2).\)
Shifting right removes the final bit. The expression \texttt{i \& 1} is one
for an odd number and zero for an even number. So the number of set
bits is the answer for the shifted value plus the removed bit.

\begin{dryrunbox}{Binary states from 0 through 5}
\begin{center}
\begin{tabular}{@{}c c c c c@{}}
\toprule
\(i\) & Binary & \(i>>1\) & Final bit & \(bits[i]\) \\
\midrule
0 & 000 & -- & -- & 0 \\
1 & 001 & 0 & 1 & \(bits[0]+1=1\) \\
2 & 010 & 1 & 0 & \(bits[1]+0=1\) \\
3 & 011 & 1 & 1 & \(bits[1]+1=2\) \\
4 & 100 & 2 & 0 & \(bits[2]+0=1\) \\
5 & 101 & 2 & 1 & \(bits[2]+1=2\) \\
\bottomrule
\end{tabular}
\end{center}
\end{dryrunbox}

\subsection{Code Explanation, Alternatives, and Edge Cases}
The output vector is filled with zeros, so the base case is already
present. The loop starts at one; \texttt{value >> 1} is always smaller and
has already been computed. A direct method could count bits separately for
every integer in \(O(n\log n)\). Another useful update rule is
\(bits[i]=bits[i\&(i-1)]+1\), which removes the lowest set bit.

For \(n=0\), the output is \texttt{[0]}. The vector is required output space,
so it should not be described as avoidable extra memory.

\begin{revisionbox}{How to explain it in an interview}
I reuse the count for the number with its final binary digit removed. Right
shift gives that smaller number, and \texttt{i \& 1} tells whether the removed
digit contributes one. Each answer is built in constant time, so the entire
range takes \(O(n)\).
\end{revisionbox}
\chaptercheckpoint
  {Each number reduces to a smaller number whose answer is known.}
  {The required output array of bit counts from zero through $n$.}
  {Recounting every binary representation independently.}
  {$O(n)$ time and $O(n)$ required output storage.}
\judgeclassnote


\patternintoc{Directed Graphs}
\problemstart{Course Schedule}{207}{Medium}{DFS cycle detection}{Directed graph}

\subsection{The Problem}
Courses are vertices. A prerequisite pair \([a,b]\) means \(b\rightarrow a\):
course \(b\) must be completed before \(a\). Return whether every course can
be finished.

\subsection{Key Observation}
A valid schedule exists exactly when the prerequisite graph has no directed
cycle. The three-state method tracks both fully visited nodes and nodes in the current DFS
path. A three-state array expresses the same idea:
\begin{itemize}
  \item \(0\): unvisited,
  \item \(1\): visiting in the current recursion path,
  \item \(2\): completely processed and safe.
\end{itemize}

\begin{invariantbox}{state 1 is exactly the active DFS path}
A course changes from unvisited to visiting when its DFS frame begins and
changes to done only after all outgoing edges are proved safe. So an
edge to state 1 is a back edge and a cycle, whereas an edge to state 2 can be
reused without more recursion.
\end{invariantbox}

\begin{visualbox}{Following the arrows returns to course 0}
\centering
\begin{tikzpicture}[>={Stealth[length=2.2mm,width=1.55mm]},font=\small,node distance=18mm]
  \node[trienode,active] (a) {0};
  \node[trienode,active,right=of a] (b) {1};
  \node[trienode,active,right=of b] (c) {2};
  \draw[flowarrow,orange] (a)--(b); \draw[flowarrow,orange] (b)--(c);
  \draw[flowarrow,orange] (c) to[bend left=35] node[edgelabel,above]{returns to active course} (a);
\end{tikzpicture}
\end{visualbox}

\subsection{Pseudocode}
\begin{lstlisting}[style=pseudocode]
hasCycle(course):
    if state[course] == visiting: return true
    if state[course] == done: return false
    state[course] = visiting
    for next in graph[course]:
        if hasCycle(next): return true
    state[course] = done
    return false
\end{lstlisting}

\subsection{C++ Implementation}
\begin{lstlisting}[style=compactcode,caption={Complete C++ solution: Course Schedule}]
#include <vector>

using namespace std;

class Solution {
    bool hasCycle(int course,
                  const vector<vector<int>>& graph,
                  vector<int>& state) {
        if (state[course] == 1) return true;
        if (state[course] == 2) return false;
        state[course] = 1;
        for (const int next : graph[course]) {
            if (hasCycle(next, graph, state)) return true;
        }
        state[course] = 2;
        return false;
    }

public:
    bool canFinish(int numCourses,
                   vector<vector<int>>& prerequisites) {
        vector<vector<int>> graph(numCourses);
        for (const auto& edge : prerequisites) graph[edge[1]].push_back(edge[0]);
        vector<int> state(numCourses, 0);
        for (int course = 0; course < numCourses; ++course) {
            if (hasCycle(course, graph, state)) return false;
        }
        return true;
    }
};
\end{lstlisting}

\subsection{Why It Works, Time, and Space}
A visiting neighbor is an ancestor in the current DFS path, so that edge closes
a directed cycle. Done nodes have already been proved cycle-free and need not
be explored again. Time and storage are \(O(V+E)\); recursion depth is
\(O(V)\) in the worst case.

\begin{revisionbox}{How to explain it in an interview}
I build the prerequisite graph and run DFS from every course. A three-state
array distinguishes an active recursion path from a completed node. Reaching
an active node proves a cycle, so the courses cannot all be finished.
\end{revisionbox}

\subsection{Why One Visited Boolean Is Not Enough}
A node seen in a different completed DFS branch is safe; a node seen again in
the current recursion path proves a cycle. A single Boolean cannot distinguish
those situations. The three states mean unvisited, visiting, and completely
processed. Only an edge to a visiting node is a back edge.

\subsection{Detailed DFS Trace}
\begin{dryrunbox}{Cycle in prerequisites \texttt{[[1,0],[0,1]]}}
\begin{center}
\begin{tabularx}{\linewidth}{@{}c c X@{}}
\toprule
Action & States \((0,1)\) & Meaning \\
\midrule
Enter course 0 & \((1,0)\) & 0 is on the active path \\
Follow \(0\to1\) & \((1,1)\) & 1 joins the active path \\
Follow \(1\to0\) & \((1,1)\) & 0 is already visiting; cycle found \\
\bottomrule
\end{tabularx}
\end{center}
For \texttt{[[1,0]]}, course 1 finishes, changes to state 2, and no back edge
is found.
\end{dryrunbox}

\begin{visualbox}{Three labels separate new, active, and finished courses}
\centering
\begin{tikzpicture}[font=\small,>={Stealth[length=2.2mm,width=1.55mm]}]
  \node[trienode,active] (c0) {0};
  \node[trienode,active,right=18mm of c0] (c1) {1};
  \draw[flowarrow,orange] (c0) to[bend left=18] (c1);
  \draw[flowarrow,orange] (c1) to[bend left=18] (c0);

  \node[right=32mm of c1,text=navy,font=\bfseries] (label) {course labels};
  \node[arraycell,active,below=5mm of label,xshift=-8mm] (s0) {1};
  \node[arraycell,active,right=5mm of s0] (s1) {1};
  \node[below=2mm of s0,font=\scriptsize] {course 0};
  \node[below=2mm of s1,font=\scriptsize] {course 1};
  \node[below=14mm of s0,anchor=north west,xshift=-2mm,
    align=left,font=\scriptsize,text=navy] {
    0 = not started\\1 = being checked now\\2 = finished and safe};
\end{tikzpicture}

\smallskip
An arrow to a course labeled 1 returns to the path being checked, so there is
a cycle.
\end{visualbox}

\begin{visualbox}{Course arrows and the path being checked}
\centering
\begin{tikzpicture}[font=\small,>={Stealth[length=2.2mm,width=1.55mm]}]
  \node[dsnode,align=left,fill=softteal] (adj) {
    \textbf{next course}\\
    \texttt{0: [1]}\\
    \texttt{1: [0]}};
  \node[draw=navy,fill=softblue,minimum width=44mm,minimum height=8mm,
    right=30mm of adj] (f0) {check course 0};
  \node[draw=navy,fill=softblue,minimum width=44mm,minimum height=8mm,
    above=0mm of f0] (f1) {check course 1};
  \node[draw=orange,fill=softorange,minimum width=44mm,minimum height=8mm,
    above=0mm of f1] (f2) {see course 0 again};
  \node[right=7mm of f2,text=orange,font=\bfseries,align=left] {working\\now};
  \draw[flowarrow,teal] (adj.east)--node[edgelabel,above,font=\scriptsize]{follow arrow} (f0.west);
  \node[below=8mm of f0,draw=orange,fill=softorange,rounded corners,
    font=\scriptsize] {course 0 is already on this path \(\Rightarrow\) cycle};
\end{tikzpicture}
\end{visualbox}

\subsection{Code Walkthrough}
The adjacency list stores edges from each prerequisite to the course unlocked
by it. The helper returns immediately for visiting and done states. Before
exploring neighbors it marks the course visiting; after every neighbor is
proved safe, it marks the course done. The outer loop is needed because
the graph may contain several disconnected components.

\subsection{Alternative, Edge Cases, and Follow-ups}
Kahn's topological-sort algorithm is an equally strong alternative: repeatedly
take zero-indegree courses and check whether all courses are removed. An empty
prerequisite list is immediately possible. A self-dependency such as
\([0,0]\) is a cycle. Duplicate edges do not change the DFS result, though
they add repeated adjacency work.

\begin{visualbox}{Two methods notice the same impossible loop}
\centering
\begin{tikzpicture}[font=\small,>={Stealth[length=2.2mm,width=1.55mm]}]
  \node[dsnode,fill=softorange] (dfs) {path view\\returns to an active course};
  \node[dsnode,right=24mm of dfs,fill=softteal] (kahn)
    {ready-list view\\no course is ready};
  \node[dsnode,below=12mm of $(dfs)!0.5!(kahn)$] (blocked)
    {the loop leaves no valid next course};
  \draw[flowarrow,orange] (dfs)--(blocked);
  \draw[flowarrow,teal] (kahn)--(blocked);
\end{tikzpicture}
\end{visualbox}

\begin{revisionbox}{Practice checkpoint}
Modify the solution to return one valid course order. Decide whether DFS
postorder or Kahn's queue makes that extension easier to explain.
\end{revisionbox}
\chaptercheckpoint
  {Prerequisites form directed dependencies that must admit an ordering.}
  {An adjacency list and three DFS states, or indegrees and a queue.}
  {Treating every previously visited node as a cycle instead of only active nodes.}
  {$O(V+E)$ time and $O(V+E)$ graph/traversal space.}
\judgeclassnote


\patternintoc{String DP}
\problemstart{Decode Ways}{91}{Medium}{Prefix dynamic programming}{String and DP}

\subsection{The Problem}
Digits map from \(1\) through \(26\) to letters. Count the ways to split a digit
string into valid codes. A zero cannot be decoded by itself.

\begin{conceptbox}{Examples}
\texttt{"12"} has two decodings: \texttt{1|2} and \texttt{12}.
\texttt{"226"} has three. \texttt{"06"} has zero because a code cannot start
with zero.
\end{conceptbox}

\subsection{DP State}
Let \(dp[i]\) be the number of decodings for the first \(i\) characters.
Start \(dp[0]=1\) for the empty prefix. At position \(i\):
\begin{itemize}
  \item if the one-digit code is \(1\) through \(9\), add \(dp[i-1]\);
  \item if the two-digit code is \(10\) through \(26\), add \(dp[i-2]\).
\end{itemize}

\begin{invariantbox}{each prefix count contains every valid ending once}
After processing length \(i\), \(dp[i]\) counts exactly the decodings of the
first \(i\) characters. Every decoding ends with either one nonzero digit or
one valid value from 10 through 26, so the previous groups are complete
and do not overlap.
\end{invariantbox}

\begin{visualbox}{For \texttt{226}, count endings \texttt{6} and \texttt{26}}
\centering
\begin{tikzpicture}[font=\small,>={Stealth[length=2.2mm,width=1.55mm]}]
  \node[arraycell,minimum width=12mm] (s0) {2};
  \node[arraycell,minimum width=12mm,right=0mm of s0] (s1) {2};
  \node[arraycell,minimum width=12mm,right=0mm of s1,active] (s2) {6};
  \node[below=2mm of s0,font=\scriptsize] {1};
  \node[below=2mm of s1,font=\scriptsize] {2};
  \node[below=2mm of s2,font=\scriptsize] {3};

  \node[dsnode,above=15mm of s1,xshift=-10mm] (one) {end with \texttt{6}\\2 earlier ways};
  \node[dsnode,above=15mm of s2,xshift=14mm,fill=softorange] (two) {end with \texttt{26}\\1 earlier way};
  \node[dsnode,right=20mm of s2,active] (ways) {total\\\(2+1=3\)};
  \draw[flowarrow,teal] (one)--(ways.north west);
  \draw[flowarrow,orange] (two)--(ways.north);
  \foreach \x/\v in {0/1,1/1,2/2,3/3} {
    \node[arraycell,minimum width=11mm] (decode\x) at ({0.1+1.25*\x},-1.55) {\v};
    \node[below=1mm of decode\x,font=\scriptsize] {\(dp[\x]\)};
  }
  \node[left=8mm of decode0,text=navy,font=\bfseries,align=right] {ways after\\each prefix};
  \node[draw=orange,very thick,fit=(decode3),inner sep=1mm] {};
\end{tikzpicture}

\smallskip
Every decoding ends with either \texttt{6} or \texttt{26}. These are separate
groups, so their counts are added.
\end{visualbox}

\begin{dryrunbox}{Decoding \texttt{"226"}}
\begin{center}
\begin{tabular}{c|cccc}
Prefix length & 0 & 1 & 2 & 3 \\ \hline
Prefix & empty & 2 & 22 & 226 \\
Ways & 1 & 1 & 2 & 3
\end{tabular}
\end{center}
\end{dryrunbox}

\subsection{Pseudocode}
\begin{lstlisting}[style=pseudocode]
if text is empty or starts with 0: return 0
dp[0] = 1
dp[1] = 1
for prefixLength from 2 through text.length:
    if last digit is 1 through 9:
        dp[prefixLength] += dp[prefixLength - 1]
    if last two digits form 10 through 26:
        dp[prefixLength] += dp[prefixLength - 2]
return dp[text.length]
\end{lstlisting}

\subsection{C++ Implementation}
\begin{lstlisting}[style=compactcode,caption={Complete C++ solution: Decode Ways}]
#include <string>
#include <vector>

using namespace std;

class Solution {
public:
    int numDecodings(string s) {
        if (s.empty() || s[0] == '0') return 0;
        const int n = static_cast<int>(s.size());
        vector<int> dp(n + 1, 0);
        dp[0] = 1;
        dp[1] = 1;

        for (int i = 2; i <= n; ++i) {
            if (s[i - 1] != '0') dp[i] += dp[i - 1];
            const int twoDigits = (s[i - 2] - '0') * 10 + (s[i - 1] - '0');
            if (twoDigits >= 10 && twoDigits <= 26) dp[i] += dp[i - 2];
        }
        return dp[n];
    }
};
\end{lstlisting}

\subsection{Why It Works, Complexity, and Common Mistakes}
Every valid decoding ends in exactly one valid one-digit or two-digit code, so
the update rule splits all possibilities without overlap. Time and space
are \(O(n)\). Carefully reject standalone zero, values above 26, and prefixes
such as \texttt{"06"}. The DP can be reduced to \(O(1)\) space because only
the previous two entries are needed.

\subsection{Brute Force and Repeated Prefixes}
A recursive solution may decode one digit, decode two digits when valid, and
recurse on the remaining suffix. For a string of many ones, both branches are
available repeatedly and the recursion grows exponentially. The same suffix
or prefix length is solved many times. Dynamic programming stores the number
of decodings for each prefix exactly once.

\subsection{Detailed Prefix Dry Run}
\begin{dryrunbox}{Decoding \texttt{"226"}}
\begin{center}
\begin{tabularx}{\linewidth}{@{}c c c X@{}}
\toprule
Prefix & One-digit choice & Two-digit choice & Ways \\
\midrule
empty & -- & -- & \(dp[0]=1\) \\
\texttt{2} & \texttt{2} is valid & -- & \(dp[1]=1\) \\
\texttt{22} & append \texttt{2}: \(dp[1]\) & append \texttt{22}: \(dp[0]\) & 2 \\
\texttt{226} & append \texttt{6}: \(dp[2]\) & append \texttt{26}: \(dp[1]\) & 3 \\
\bottomrule
\end{tabularx}
\end{center}
The three decodings are \texttt{2|2|6}, \texttt{22|6}, and
\texttt{2|26}.
\end{dryrunbox}

\subsection{Code Walkthrough}
The first-character guard rejects an empty input or leading zero. The table
uses prefix lengths, so \texttt{dp[i]} describes the first \(i\) characters.
If the last character is nonzero, every decoding of the previous prefix can
append it. If the last two characters form 10 through 26, every decoding two
positions back can append that pair.

\subsection{Important Zero Cases and Follow-ups}
\texttt{"10"} and \texttt{"20"} each have one decoding. \texttt{"30"},
\texttt{"06"}, and \texttt{"100"} are invalid. The pair test must require at
least 10, not merely at most 26. Since only two earlier states are read, the
table can be compressed to two integers for \(O(1)\) extra space.

\begin{dryrunbox}{Zero-heavy strings that expose invalid transitions}
\begin{center}
\begin{tabularx}{\linewidth}{@{}c c X@{}}
\toprule
String & Ways & Reason \\
\midrule
\texttt{0} & 0 & zero cannot stand alone \\
\texttt{10} & 1 & only the pair 10 is valid \\
\texttt{101} & 1 & \texttt{10|1}; zero blocks the one-digit transition \\
\texttt{100} & 0 & the final zero cannot pair with zero \\
\texttt{226} & 3 & \texttt{2|2|6}, \texttt{22|6}, and \texttt{2|26} \\
\bottomrule
\end{tabularx}
\end{center}
\end{dryrunbox}

\begin{revisionbox}{How to explain it simply}
Every decoding ends with either one valid digit or one valid two-digit code.
Those two groups do not overlap, so I add the matching prefix counts. The
DP scans once in \(O(n)\) time and explicitly handles zeros.
\end{revisionbox}
\chaptercheckpoint
  {A prefix can end with one-character or two-character codes.}
  {The number of valid decodings for each processed prefix.}
  {Allowing zero alone or accepting a two-digit number outside 10--26.}
  {$O(n)$ time and $O(n)$ space, reducible to $O(1)$.}
\judgeclassnote


\patternintoc{Linked Lists}
\problemstart{Linked List Cycle}{141}{Easy}{Floyd's two-pointer cycle detection}{Linked list}

\subsection{The Problem}
Return whether repeatedly following \texttt{next} can revisit a node. The
platform's cycle position is not passed to the function.

\subsection{Why Slow and Fast Pointers Meet}
A hash set can store visited pointers in \(O(n)\) space. Floyd's algorithm uses
two pointers: slow advances one edge and fast advances two. If there is no
cycle, fast reaches null. Inside a cycle, fast gains one node per iteration on
slow and must eventually meet it.

\begin{invariantbox}{null proves an end; pointer equality proves repetition}
After \(k\) rounds, slow has moved \(k\) edges and fast has moved \(2k\). If
fast reaches null, the list is finite. If both pointers reference the same
node, different traversal distances have reached one object, which can happen
only after entering a cycle.
\end{invariantbox}

\begin{visualbox}{Fast catches slow when the list loops}
\centering
\begin{tikzpicture}[>={Stealth[length=2.2mm,width=1.55mm]},font=\small,node distance=14mm]
  \node[trienode] (a) {3}; \node[trienode,right=of a] (b) {2};
  \node[trienode,right=of b] (c) {0}; \node[trienode,right=of c] (d) {-4};
  \draw[flowarrow] (a)--(b)--(c)--(d);
  \draw[flowarrow,orange] (d) to[bend right=45] node[edgelabel,below]{loop back} (b);
  \node[above=7mm of b,text=teal] {slow};
  \node[above=7mm of c,text=orange] {fast};
\end{tikzpicture}
\end{visualbox}

\subsection{Pseudocode}
\begin{lstlisting}[style=pseudocode]
slow = head
fast = head
while fast exists and fast.next exists:
    slow = slow.next
    fast = fast.next.next
    if slow equals fast: return true
return false
\end{lstlisting}

\subsection{C++ Implementation}
\begin{lstlisting}[style=compactcode,caption={Complete C++ solution: Linked List Cycle}]
using namespace std;

class Solution {
public:
    bool hasCycle(ListNode* head) {
        ListNode* slow = head;
        ListNode* fast = head;
        while (fast != nullptr && fast->next != nullptr) {
            slow = slow->next;
            fast = fast->next->next;
            if (slow == fast) return true;
        }
        return false;
    }
};
\end{lstlisting}

\subsection{Time and Space}
Both pointers make \(O(n)\) total progress before fast reaches null or the
pointers meet. Extra space is \(O(1)\). Check both \texttt{fast} and
\texttt{fast->next} before advancing two steps.

\begin{revisionbox}{Pattern clue}
Use slow and fast pointers when repeated next-state transitions may enter a
cycle and the problem asks for constant extra space.
\end{revisionbox}

\subsection{From a Hash Set to Constant Space}
The simplest solution stores every visited node pointer in a hash set. Before
advancing, it checks whether the current pointer was seen before. This is
linear time and linear space. Floyd's method removes the set by using two
walkers with different speeds.

If a cycle exists, slow and fast eventually enter it. After that, fast gains
one node on slow per iteration, measured around the cycle. The distance
modulo the cycle length must eventually become zero, so the pointers meet.

\subsection{Pointer Dry Run}
\begin{dryrunbox}{List \(3\to2\to0\to-4\), with \(-4\to2\)}
\begin{center}
\begin{tabular}{@{}c c c@{}}
\toprule
Iteration & Slow & Fast \\
\midrule
start & 3 & 3 \\
1 & 2 & 0 \\
2 & 0 & 2 \\
3 & -4 & -4 \\
\bottomrule
\end{tabular}
\end{center}
The equal pointers refer to the same node object, not merely equal stored
values. So a repeated value in an acyclic list does not create a false
cycle.
\end{dryrunbox}

\begin{visualbox}{Without a loop, fast reaches the end}
\centering
\begin{tikzpicture}[font=\small,>={Stealth[length=2.2mm,width=1.55mm]}]
  \node[dsnode] (a) {1}; \node[dsnode,right=10mm of a] (b) {2};
  \node[dsnode,right=10mm of b] (c) {3}; \node[dsnode,right=10mm of c] (d) {4};
  \node[right=10mm of d] (null) {end};
  \draw[flowarrow] (a)--(b)--(c)--(d)--(null);
  \node[above=7mm of b,text=teal,font=\bfseries] {slow};
  \node[above=7mm of d,text=orange,font=\bfseries] {fast};
  \node[below=8mm of d,text=navy] {fast reaches the end \(\Rightarrow\) no loop};
\end{tikzpicture}
\end{visualbox}

\subsection{Code Explanation and Safety}
Both pointers begin at the head. The loop verifies that fast and
\texttt{fast->next} exist before moving fast two edges. The comparison occurs
after movement; if a cycle exists, the pointers eventually become identical.
If fast reaches null, the list has a finite end.

\subsection{Edge Cases and Interview Follow-ups}
Test an empty list, one node without a loop, one node pointing to itself, a
two-node cycle, and a long non-cyclic prefix before a short cycle. A common
follow-up asks for the cycle's entry node: after the first meeting, reset one
pointer to head and advance both one step at a time; their next meeting is the
entry.

\begin{revisionbox}{Practice checkpoint}
Explain why comparing \texttt{slow->val == fast->val} is incorrect. Then
derive how resetting one pointer after the meeting locates the cycle entry.
\end{revisionbox}
\chaptercheckpoint
  {Repeated next-pointer transitions may enter a cycle.}
  {Slow and fast node pointers.}
  {Comparing values instead of pointer identity or dereferencing past null.}
  {$O(n)$ time and $O(1)$ extra space.}
\judgeclassnote


\patternintoc{String Design}
\problemstart{Encode and Decode Strings}{271}{Medium}{Length-prefixed encoding}{Strings}

\subsection{The Problem}
Design two operations: encode a vector of arbitrary strings into one string,
then decode it back exactly. Strings may be empty and may contain delimiter
characters.

\subsection{Why a Delimiter Alone Fails}
Joining with \texttt{\#} is ambiguous when a payload itself contains
\texttt{\#}. Prefix every string with its decimal length and a delimiter: \(\boxed{\texttt{length\#payload}}.\)
The decoder reads the length first, so it knows exactly how many following
characters belong to the payload.

\begin{invariantbox}{the cursor always begins at the next frame length}
After decoding a frame, the cursor has advanced past its decimal length, the
separator, and exactly the announced payload characters. Separators inside a
payload are never interpreted, so the next cursor position is clear.
\end{invariantbox}

\begin{visualbox}{Each stored string says its own length}
\centering
\begin{tikzpicture}[font=\small]
  \node[dsnode] (a) {\texttt{4\#lint}};
  \node[dsnode,right=4mm of a] (b) {\texttt{4\#code}};
  \node[dsnode,right=4mm of b] (c) {\texttt{0\#}};
  \node[dsnode,right=4mm of c] (d) {\texttt{3\#a\#b}};
  \node[below=7mm of $(a)!0.5!(d)$] {the number says how many characters come next};
\end{tikzpicture}
\end{visualbox}

\subsection{Pseudocode}
\begin{lstlisting}[style=pseudocode]
encode(strings): append length + '#' + string for every string
decode(data):
    i = 0
    while i < data.length:
        find next '#'
        parse length before it
        copy exactly length characters after it
        advance i to the next frame
\end{lstlisting}

\subsection{C++ Implementation}
\begin{lstlisting}[style=compactcode,caption={Complete C++ solution: Encode and Decode Strings}]
#include <string>
#include <vector>

using namespace std;

class Codec {
public:
    string encode(const vector<string>& strings) {
        string encoded;
        for (const string& value : strings) {
            encoded += to_string(value.size());
            encoded.push_back('#');
            encoded += value;
        }
        return encoded;
    }

    vector<string> decode(const string& encoded) {
        vector<string> result;
        size_t index = 0;
        while (index < encoded.size()) {
            const size_t delimiter = encoded.find('#', index);
            const size_t length =
                static_cast<size_t>(stoull(encoded.substr(index, delimiter-index)));
            const size_t start = delimiter + 1;
            result.push_back(encoded.substr(start, length));
            index = start + length;
        }
        return result;
    }
};
\end{lstlisting}

\subsection{Why It Works, Complexity, and Edge Cases}
Every frame has one clear numeric prefix. The decoder consumes exactly
the announced payload length, even when the payload contains \texttt{\#} or is
empty. Encoding and decoding both take \(O(N)\) time for total character count
\(N\), with \(O(N)\) returned storage. Test empty vectors, empty strings,
delimiter characters, and multi-digit lengths.

\begin{revisionbox}{How to explain it in an interview}
I avoid delimiter ambiguity by storing the payload length before each string.
The decoder finds the separator, parses the length, copies exactly that many
characters, and continues at the next frame.
\end{revisionbox}

\subsection{Why Simple Joining Is Not Reversible}
Suppose the separator is \texttt{\#}. Encoding
\texttt{["ab","c\#d"]} as \texttt{ab\#c\#d} loses the information that the
second separator belongs inside a payload. Escaping separators is possible,
but then the escape character itself must also be escaped. A length prefix is
simpler because payload contents no longer affect parsing.

\subsection{Frame-by-Frame Dry Run}
\begin{dryrunbox}{Encoding \texttt{["lint","code","","a\#b"]}}
\begin{center}
\begin{tabular}{@{}c c c@{}}
\toprule
String & Frame & Decoder action \\
\midrule
\texttt{lint} & \texttt{4\#lint} & read 4, then copy four characters \\
\texttt{code} & \texttt{4\#code} & read 4, then copy four characters \\
empty & \texttt{0\#} & read 0 and append an empty string \\
\texttt{a\#b} & \texttt{3\#a\#b} & copy three characters, including \texttt{\#} \\
\bottomrule
\end{tabular}
\end{center}
\end{dryrunbox}

\subsection{Code Walkthrough}
Encoding appends the decimal length, one separator, and the untouched payload.
Decoding searches only for the separator that terminates the numeric length.
It parses that prefix, calculates the payload start, copies exactly the stated
number of characters, and advances directly to the next frame. The platform
supplies strings produced by the matching encoder, so invalid input handling
is outside this problem's contract.

\subsection{Edge Cases, Complexity Details, and Variations}
Distinguish an empty vector from a vector containing one empty string: their
encodings are respectively empty and \texttt{0\#}. Multi-digit lengths must
be parsed completely. Total work is linear in the encoded data size, although
repeated string growth may allocate; reserving capacity is a practical
optimization when total size is known.

\medskip
\noindent\textcolor{teal}{\textbf{Practice checkpoint.}}
Decode \texttt{5\#hello0\#3\#a\#b} by hand. At each step record the separator
position, parsed length, payload interval, and next frame index.
\chaptercheckpoint
  {Arbitrary strings must be combined and separated without ambiguity.}
  {Length-prefixed frames and a decoder cursor.}
  {Splitting on a delimiter that may legally occur inside the payload.}
  {$O(N)$ encoding/decoding time and $O(N)$ returned storage.}
\judgeclassnote